\reading{IM-6}{VaR and Risk Budgeting in Investment Management}{STRIKE FIRST}{4}

\srcnote{Philippe Jorion, \textit{Value-at-Risk: The New Benchmark for Managing
Financial Risk}, 3rd Edition (New York, NY: McGraw-Hill, 2007), Chapter 17.
Printed as Chapter 6 in the GARP source volume.}

\begin{imkeybox}[title={The eight objectives in this reading}]
\begin{enumerate}[label=\textbf{\alph*.}]
\item Define risk budgeting.
\item Describe the impact of horizon, turnover, and leverage on the risk management process in the investment management industry.
\item Describe the investment process of large investors such as pension funds.
\item Describe the risk management challenges associated with investments in hedge funds.
\item Describe and compare the following types of risk: absolute risk, relative risk, policy-mix risk, active management risk, funding risk, and sponsor risk.
\item Explain the use of VaR to monitor risk.
\item Explain how VaR can be used in the development of investment guidelines and for improving the investment process.
\item Describe the risk budgeting process and calculate risk budgets across asset classes and active managers.
\end{enumerate}
\end{imkeybox}

% ------------------------------------------------------------------
\lo{IM-6 a}{Define risk budgeting.}

\begin{imdefbox}[title={Risk budgeting}]
A \term{top-down} approach where a risk manager assigns a risk limit to the
entire portfolio and distributes this risk among various investments based on
their targeted risk levels. It differs from market value allocation in that it
focuses on \term{risk}, not just investment value.
\end{imdefbox}

% ------------------------------------------------------------------
\lo{IM-6 b}{Describe the impact of horizon, turnover, and leverage on the risk management process in the investment management industry.}

\renewcommand{\arraystretch}{1.25}
\noindent\begin{tabularx}{\textwidth}{@{}>{\raggedright\arraybackslash}p{2.2cm}XX@{}}
\toprule
& \textbf{\sffamily\small Sell-side (banks)} & \textbf{\sffamily\small Buy-side (pension funds, institutions)} \\
\midrule
Horizon & Short. Positions change frequently. & Long-term holdings. \\
Turnover & High --- rapid trading. & Low. \\
Leverage & High. & Lesser leverage. \\
Consequence & Needs a \term{responsive, dynamic} measure because positions change constantly. & VaR must adapt to long horizons and to increased complexity in global investments. \\
\bottomrule
\end{tabularx}
\renewcommand{\arraystretch}{1}
\figcap{The same tool, two different jobs. The sell-side uses VaR because its
book turns over faster than any static limit could track; the buy-side uses it
because its holdings are long-dated and complex enough that intuition about
aggregate risk fails. Horizon, turnover and leverage all point the same way for
each side, which is why the two industries ended up with different VaR practices
from the same measure.}

% ------------------------------------------------------------------
\lo{IM-6 c}{Describe the investment process of large investors such as pension funds.}

\begin{enumerate}
\item In the first step, a consultant provides a \term{strategic, long-term
asset-allocation study} that balances expected return against risk. This study
determines the amounts to be invested in various asset classes.
\item In the second step, the fund may \term{delegate the actual management} of
funds to a stable of active managers. These managers are reviewed periodically
for performance relative to their benchmark.
\end{enumerate}

\begin{imkeybox}
The two steps map exactly onto the two risks in IM-6 e. Step one creates
\term{policy-mix risk}; step two creates \term{active management risk}.
\end{imkeybox}

% ------------------------------------------------------------------
\lo{IM-6 d}{Describe the risk management challenges associated with investments in hedge funds.}

\begin{itemize}
\item Hedge funds are \term{heterogeneous}, and so pose special risk measurement
problems.
\item Most hedge funds have \term{leverage}.
\item Hedge funds prefer investing in \term{illiquid assets}. In this case, risk
measures based on monthly returns give a misleading picture of risk.
\end{itemize}

\begin{imtrapbox}[breakable=false,title={Illiquidity biases both numbers the same way}]
Illiquidity causes two distinct biases, and both flatter the fund:
\begin{enumerate}
\item Correlation with other asset classes is \term{artificially lowered}, giving
the appearance of low \emph{systematic} risk.
\item Volatility is \term{artificially lowered}, giving the appearance of low
\emph{total} risk.
\end{enumerate}
Both are downward. A question that has one bias raising a risk measure has
reversed it.
\end{imtrapbox}

% ------------------------------------------------------------------
\lo{IM-6 e}{Describe and compare the following types of risk: absolute risk, relative risk, policy-mix risk, active management risk, funding risk, and sponsor risk.}

\figwrap{%
\begin{tikzpicture}[
  hd/.style={draw=FRMNavy,fill=FRMNavy,rounded corners=1pt,inner sep=4pt,
             font=\scriptsize\sffamily\bfseries,text=white,align=center,
             text width=36mm},
  bx/.style={draw=FRMRule,fill=white,rounded corners=1pt,inner sep=4pt,
             font=\scriptsize\sffamily,align=center,text width=36mm},
  ar/.style={-{Stealth[length=4pt]},FRMGold,line width=0.9pt}]
\node[hd] (ab) at (0,0) {Absolute risk};
\node[bx,below=4mm of ab] (ab1) {dollar loss measured\\ \textbf{relative to zero}};
\draw[ar] (ab) -- (ab1);
\node[bx,below=6mm of ab1] (pm) {\textbf{Policy-mix risk}\\ loss owing to the policy\\ mix the fund selected};
\node[bx,below=3mm of pm] (am) {\textbf{Active management risk}\\ loss owing to deviations\\ from the policy mix};
\draw[ar] (ab1) -- node[right,font=\scriptsize\sffamily,text=black!60,
  fill=white,inner sep=1.5pt]{splits into} (pm);

\node[hd] (re) at (5.2,0) {Relative risk};
\node[bx,below=4mm of re] (re1) {dollar loss measured\\ \textbf{relative to a benchmark}};
\draw[ar] (re) -- (re1);

\node[hd] (fu) at (10.4,0) {Funding risk};
\node[bx,below=4mm of fu] (fu1) {assets insufficient\\ to fund liabilities};
\node[bx,below=3mm of fu1] (fu2) {measured as the shortfall\\ in \textbf{surplus at risk}};
\draw[ar] (fu) -- (fu1); \draw[ar] (fu1) -- (fu2);
\node[hd,below=6mm of fu2] (sp) {Sponsor risk};
\node[bx,below=3mm of sp] (sp1) {\textbf{Cash-flow risk:} year-to-year\\ fluctuations in contributions};
\node[bx,below=3mm of sp1] (sp2) {\textbf{Economic risk:} variation in\\ the sponsor's total earnings};
\draw[ar] (sp) -- (sp1); \draw[ar] (sp1) -- (sp2);
\end{tikzpicture}}%
{Six risks and the three questions that separate them. \emph{Against what} is the
loss measured --- zero or a benchmark? That splits absolute from relative.
\emph{Who made the decision} that caused it --- the consultant who set the policy
mix, or the manager who deviated from it? That splits policy-mix from active
management risk, and both live inside absolute risk. \emph{Whose problem is it} ---
the fund's, or the firm that stands behind the fund? That splits funding risk
from sponsor risk. Note that policy-mix and active management risk are
\emph{components} of absolute risk, not alternatives to it.}

\subsection{Funding risk and surplus at risk}

Surplus $S$ is the difference between the value of the assets $A$ and the
liabilities $L$, so $\Delta S = \Delta A - \Delta L$. Funding risk should be
measured as the potential shortfall in surplus over the horizon, sometimes called
\term{surplus at risk (SAR)}.

\begin{imfmlbox}
\[ \text{Expected surplus} = A\,(1 + R_A) - L\,(1 + R_L) \]
\vk{$R_A$ is the return on assets and $R_L$ the growth rate of liabilities. The
source notes print this without parentheses, as $A \times 1 + R_A - L \times 1 +
R_L$, which evaluates to something else entirely.}
\end{imfmlbox}

\begin{imtrapbox}[breakable=false,title={Sponsor risk is measured twice over}]
The sponsor is the \emph{owner} of the fund, ultimately responsible for it.
Sponsor risk is measured both by the movements in the assets --- or even the
surplus --- \emph{and} by the ultimate effect on the economic value of the firm.
It is the only risk in this list that reaches outside the fund.
\end{imtrapbox}

% ------------------------------------------------------------------
\lo{IM-6 f}{Explain the use of VaR to monitor risk.}

\begin{enumerate}
\item VaR sets risk budgets for managers and identifies deviations.
\item It can detect \term{unauthorized trades}, \term{portfolio drift}, or
increased risk-taking.
\item VaR helps investigate reasons for risk changes, such as manager behaviour,
correlated bets, or market volatility.
\item Advanced VaR tools can predict how individual investments impact overall
risk.
\item While VaR has limitations, it is a valuable tool for proactive risk
management and better decision-making.
\item Centralized risk management with global custodians can further improve risk
monitoring.
\end{enumerate}

% ------------------------------------------------------------------
\lo{IM-6 g}{Explain how VaR can be used in the development of investment guidelines and for improving the investment process.}

\begin{itemize}
\item \term{For guidelines:} VaR replaces vague, ad-hoc rules with clear risk
limits that consider factors like leverage and diversification. This helps ensure
managers stay within acceptable risk boundaries.
\item \term{For the investment process:} VaR helps with strategic asset allocation
by using similar techniques to those used in portfolio risk assessment. It allows
for measuring the impact of portfolio adjustments, guiding managers towards better
risk-return trade-offs.
\item VaR is also valuable \term{at the trading level}, where it helps assess how
individual trades affect the entire portfolio's risk profile.
\end{itemize}

% ------------------------------------------------------------------
\lo{IM-6 h}{Describe the risk budgeting process and calculate risk budgets across asset classes and active managers.}

Risk budgeting is the process of allocating and managing risks to different
aspects of the investment process using a \term{top-down} approach. The budget
helps the organization stay on course with respect to its risk plan. There are
two types:

\begin{itemize}
\item Budgeting \term{across asset classes};
\item Budgeting \term{across asset managers}.
\end{itemize}

\subsection{Budgeting across asset classes}

\begin{imexbox}[title={Which addition keeps the manager inside the budget?}]
A manager has a portfolio with only one position: a \$500 million investment in
W. He is considering adding a \$500 million position X or Y. The current
volatility of W is 10\%. He wants to limit portfolio VaR to \$200 million at 99\%
confidence. Position X has a return volatility of 9\% and a correlation with W of
0.7. Position Y has a return volatility of 12\% and a correlation with W of zero.
Which addition keeps him within his risk budget?

\medskip
Currently, the VaR of the portfolio with only W is
\[ \text{VaR}_W = (2.33)(10\%)(\$500\text{ million}) = \$116.5\text{ million} \]
Adding X, the return volatility of the portfolio becomes
\begin{align*}
\sigma_{W+X} &= \sqrt{(0.5^2)(10\%)^2 + (0.5^2)(9\%)^2 + 2(0.5)(0.5)(0.7)(10\%)(9\%)} = 8.76\% \\
\text{VaR}_{W+X} &= (2.33)(8.76\%)(\$1{,}000\text{ million}) = \$204\text{ million}
\end{align*}
Adding Y instead, with zero correlation the cross term vanishes:
\begin{align*}
\sigma_{W+Y} &= \sqrt{(0.5^2)(10\%)^2 + (0.5^2)(12\%)^2} = 7.81\% \\
\text{VaR}_{W+Y} &= (2.33)(7.81\%)(\$1{,}000\text{ million}) = \$182\text{ million}
\end{align*}
\textbf{Y keeps the total portfolio within the risk budget; X does not.}
\end{imexbox}

\begin{imtrapbox}[breakable=false,title={The riskier asset is the safer addition}]
Y has the \emph{higher} standalone volatility --- 12\% against X's 9\% --- and is
nevertheless the one that fits the budget, because its correlation with W is
zero. Standalone volatility does not decide this question; the contribution to
portfolio volatility does. Note too that the portfolio doubles to \$1,000 million
after either addition: using \$500 million in the second step is the standard
intermediate-result slip here.
\end{imtrapbox}

\subsection{Budgeting across active managers}

The traditional method for evaluating active managers is by measuring their
tracking error and using it to derive the information ratio.

\begin{imfmlbox}
\[ \text{IR}_i = \frac{\text{expected tracking error of the manager}}
                      {\text{volatility of the manager's tracking error}}
   = \frac{\mu_i}{\omega_i}
   \qquad\qquad
   \text{IR}_p = \frac{\mu_p}{\omega_p} \]
\vk{$\mu$ is the expected tracking error and $\omega$ its volatility, also written
TEV.}
\end{imfmlbox}

\begin{imtrapbox}[breakable=false,title={``Tracking error'' means two different things}]
GARP's own wording here defines tracking error (TE) as \emph{the active return
minus that of the benchmark} --- that is, TE is the \term{return difference
itself}, a series with a mean $\mu$ and a volatility $\omega$. In common usage,
and in IM-3 b of this volume, ``tracking error'' means $\omega$, the
\term{volatility} of that difference. Both readings are in the curriculum. The
information ratio is $\mu/\omega$ under either, so read which quantity a stem is
handing you: a \emph{level} or a \emph{standard deviation}. This is a precision
point, not an error in the notes.
\end{imtrapbox}

\begin{imfmlbox}
\[ \text{weight of portfolio managed by manager } i =
   \frac{\text{IR}_i \times (\text{portfolio's tracking error})}
        {\text{IR}_p \times (\text{manager's tracking error})} \]
\end{imfmlbox}

\gapbadge
\srcnote{The notes give the allocation formula and stop. GARP's worked case
below (Chapter 6, Table 6.4) is what shows where the numbers come from, and it
carries the result that makes risk budgeting worth doing.}

\begin{imgapbox}[title={What the objective asks for}]
The objective's verb is \emph{calculate}. Without a worked case the formula is
unusable, because the portfolio's own information ratio $\text{IR}_p$ on the
right-hand side is not given --- it is an output.
\end{imgapbox}

\begin{imexbox}[title={Allocating \$525 million across two active managers}]
A fund wants to allocate \$525 million to a pool of active managers so as to
maximise the information ratio of the fund, subject to an overall TEV of 4\%.
Each manager has a TEV of 6\%. Their information ratios are 0.60 and 0.40. To
achieve an exact TEV of 4\%, some residual investment is held in the benchmark,
which has a TEV of zero.

\medskip
Since risk budgets are proportional to information ratios, set
$x_i \omega_i \propto \text{IR}_i$. With $\omega_i = 6\%$ for both:
\[ \sqrt{(x_1\omega)^2 + (x_2\omega)^2} = 4\%
   \;\Longrightarrow\;
   x_i = \frac{0.04}{0.06\sqrt{0.60^2 + 0.40^2}} \times \text{IR}_i \]
\begin{center}\small
\begin{tabular}{@{}lrrrr@{}}
\toprule
& IR & Weight & Allocation & Risk budget (99\%) \\
\midrule
Manager 1 & 0.60 & 55\% & \$291m & \$40.6m \\
Manager 2 & 0.40 & 37\% & \$194m & \$27.1m \\
Benchmark & --- & \phantom{0}8\% & \$40m & \$0.0m \\
\midrule
Total & & 100\% & \$525m & \$48.9m \\
\bottomrule
\end{tabular}
\end{center}
\medskip
Expected excess return is $0.555(3.6\%) + 0.370(2.4\%) = 2.9\%$, where each
manager's expected TE is $\mu_i = \text{IR}_i \times \omega_i$. Against a 4\% TEV
that is a portfolio information ratio of
\[ \text{IR}_p = \frac{2.9\%}{4.0\%} = \mathbf{0.72} \]
\end{imexbox}

\begin{imkeybox}[title={The result worth remembering}]
\term{The portfolio's information ratio of 0.72 is higher than either manager's}
--- higher than the 0.60 of the better one. That is not a rounding artefact: it
is the diversification benefit from the managers' deviations being independent of
each other. Combining a 0.60 manager with a 0.40 manager produces something
better than the 0.60 manager alone.
\end{imkeybox}

\figwrap{%
\begin{tikzpicture}
\begin{axis}[frmaxis, name=leftax, scale only axis,
  width=4.6cm, height=4.0cm,
  title={Information ratios},
  ybar, bar width=20pt,
  symbolic x coords={Manager 1, Manager 2},
  xtick=data, x tick label style={font=\scriptsize\sffamily},
  enlarge x limits=0.45,
  ymin=0, ymax=0.75, ytick={0,0.2,0.4,0.6},
]
\addplot[draw=FRMNavy,fill=PaleNavy,area legend] coordinates
  {(Manager 1,0.60) (Manager 2,0.40)};
\node[font=\scriptsize\sffamily,color=black!70,anchor=south]
  at (axis cs:Manager 1,0.61) {$0.60$};
\node[font=\scriptsize\sffamily,color=black!70,anchor=south]
  at (axis cs:Manager 2,0.41) {$0.40$};
\end{axis}
\begin{axis}[frmaxis, name=rightax, at={(leftax.right of south east)}, xshift=2.0cm,
  scale only axis, width=4.6cm, height=4.0cm,
  title={Risk budgets (\$m)},
  ybar, bar width=20pt,
  symbolic x coords={Manager 1, Manager 2},
  xtick=data, x tick label style={font=\scriptsize\sffamily},
  enlarge x limits=0.45,
  ymin=0, ymax=50, ytick={0,10,20,30,40},
]
\addplot[draw=FRMGold,fill=PaleGold,area legend] coordinates
  {(Manager 1,40.6) (Manager 2,27.1)};
\node[font=\scriptsize\sffamily,color=black!70,anchor=south]
  at (axis cs:Manager 1,41) {$40.6$};
\node[font=\scriptsize\sffamily,color=black!70,anchor=south]
  at (axis cs:Manager 2,27.6) {$27.1$};
\end{axis}
\end{tikzpicture}}%
{Skill in, risk out, in the same proportion. The information ratios stand at
$0.60$ and $0.40$, a ratio of $1.50$; the risk budgets come out at \$40.6m and
\$27.1m, a ratio of $1.50$. \term{Relative risk budgets are proportional to
information ratios} --- that single sentence is the content of the optimisation,
and it is why the better manager is given more \emph{risk}, not merely more
money.}

\begin{imtrapbox}[breakable=false,title={Consolidated trap summary --- IM-6}]
\begin{itemize}
\item \term{Role, IM-6 c/e.} The consultant's asset-allocation study creates
\term{policy-mix risk}; the delegated managers create \term{active management
risk}. Both sit inside absolute risk.
\item \term{Definition, IM-6 e.} Absolute risk is measured against \emph{zero};
relative risk against a \emph{benchmark}. Everything else in the list follows
from which of the two you started with.
\item \term{Polarity, IM-6 d.} Illiquidity biases correlation \emph{down} and
volatility \emph{down}. Both understate risk; neither overstates it.
\item \term{Formula, IM-6 e.} Expected surplus is $A(1+R_A) - L(1+R_L)$. The
parentheses are load-bearing.
\item \term{Intermediate result, IM-6 h.} After adding a second \$500 million
position the portfolio is \$1,000 million. Using the original \$500 million halves
the answer and lands on a plausible figure.
\item \term{Sibling, IM-6 h.} Standalone volatility does not decide which addition
fits a risk budget. The 12\% asset with zero correlation beats the 9\% asset with
0.7 correlation.
\item \term{Definition, IM-6 h.} Risk budgets are proportional to information
ratios, not to assets under management or to expected return.
\item \term{Scope, IM-6 h.} The combined information ratio can exceed every
individual manager's, because their deviations are assumed independent.
\end{itemize}
\end{imtrapbox}

% ==================================================================
